\section{Earlier residual forecasting readouts}\label{app:earlier-decoders}

The completed residual-forecast studies read out future observations
through bounded features and retained baselines. This appendix preserves
their transformations, fitting rules, and generic ridge bounds.
Section~\ref{sec:bounded-decoder} specifies the direct affine readouts
for the G/P/K/R representation comparison, which separately transform
membership and position blocks. Neither architecture estimates or
iterates a transition matrix. Bounded decoder inputs are distinguished
from the raw observable spaces analyzed in
Section~\ref{sec:reconstruction-theory}.

\subsection{Joint common block and separate departure block}


Each feature block is centered and scaled coordinatewise using A alone;
feature standard deviations are floored at $10^{-6}$. A tilde denotes
this coordinatewise standardization. Apply the radial map
$T_R$ in Equation~\eqref{eq:mechanism-main-radial}, with $R=\sqrt p$
for a block of width $p$.
This radial map has norm at most $R$ and is nonexpansive with respect to
its argument. If the common and departure block dimensions are $p_0$
and $p_d$, respectively, set $R_0=\sqrt{p_0}$ and $R_d=\sqrt{p_d}$.
The bounded feature vectors are
\[
 \chi_K(z)=\operatorname{concat}\bigl(
       T_{R_0}(\widetilde\psi_0(z)),T_{R_d}(\widetilde d_K(z))\bigr),
 \qquad
 \chi_R(z)=\operatorname{concat}\bigl(
       T_{R_0}(\widetilde\psi_0(z)),T_{R_d}(\widetilde d_R(z))\bigr).
\]
The two blocks are transformed separately and are not standardized again.
Consequently, appending a departure does not rescale the common features.
For $J$ local models with stored ranks $r_j$, the dimensions are
$p_0=J+\sum_j r_j$ and $p_d=Jd$. The earlier dimension experiment requests
partitions with 2, 4, and 8 models and at most 24 references per model.
All fitted banks contain $J=14$ models, so $p_d=168$ for its
12-dimensional histories, whereas $p_0$ varies with the local ranks.
Each candidate geometry is fitted without targets on A and remains fixed
during selection and evaluation. All bounds use its actual dimensions.

\subsection{Fitting and calibrating the correction}

In the dimension study, let $f_{0,h}$ be the raw-history ridge baseline
for horizon $h$. Given
an unstandardized history $x$, let $\psi(x)$ denote a generic chosen
bounded feature vector, such as $\chi_K$ or $\chi_R$, after applying its
A-fitted input scaler. The forecast is
\begin{equation}
 \widehat y_{t+h}=f_{0,h}(x_t)+\gamma_h q_h(x_t),\qquad
 q_h(x)=b_h+\psi(x)^\top w_h.
 \label{eq:mechanism-main-forecast}
\end{equation}
Here $b_h\in\R$ is a free intercept, $w_h$ is the coefficient vector,
and $\gamma_h$ controls the correction strength. For $n_A$ fitting
histories, the residual targets are
$e_{ih}=y_{i,h}-f_{0,h}(x_i)$, where $y_{i,h}$ is the observed future
at horizon $h$ for fitting history $i$. The correction minimizes
\[
 \frac1{n_A}\sum_{i=1}^{n_A}
       \{e_{ih}-b_h-\psi(x_i)^\top w_h\}^2
       +\alpha_h\|w_h\|_2^2.
\]
When $\psi$ contains both common and departure blocks, their coefficients
are fitted jointly. Sharing the raw-history baseline and the common
feature construction therefore does not fix the fitted common
contribution across departure methods.
Set B selects $\alpha_h\in\{10^{-4},10^{-3},10^{-2},10^{-1},1\}$
by forecast error separately for each horizon and candidate cap.
For each geometric procedure $m$, let $\widehat y^{(c)}_{m,h}$ be its
B prediction after this penalty selection, with $\gamma_h=1$. One cap
is selected across horizons by
\begin{equation}
 \widehat c_m=\mathop{\arg\min}_{c\in\{1,2,4,8,12,20\}}
 \frac1{|\mathcal H|}\sum_{h\in\mathcal H}
 \frac{\MSE_B(Y_h,\widehat y^{(c)}_{m,h})}
 {\max\{\Var_A(Y_h),10^{-12}\}}.
 \label{eq:dimension-cap-selection}
\end{equation}
Here $\Var_A$ is the empirical variance of the retained A targets with
divisor $n_A$; exact ties select the smallest cap. The common, bounded
KAHM, bounded RBF, and two unbounded-departure procedures each make
their own choice. These choices are fixed before accessing C.
Set C then selects $\gamma_h\in\{0,1/4,1/2,1\}$ separately by
horizon. Neither selection step refits the geometry or decoder, and
set D is reserved for final evaluation. The cap objective uses
uncalibrated B forecasts, so C does not enter dimension selection.

\subsection{A uniform bound including the intercept}

Define the fitting residual mean, centered root mean square, and feature
mean as
\[
 \bar e_h=\frac1{n_A}\sum_i e_{ih},\qquad
 \sigma_{e,h}=\left[\frac1{n_A}\sum_i(e_{ih}-\bar e_h)^2\right]^{1/2},
 \qquad
 \bar\psi=\frac1{n_A}\sum_i\psi(x_i).
\]
\begin{proposition}[Centered ridge correction bound]
\label{prop:mechanism-correction-bound}
For the mean-squared-error objective above with $n_A\geq1$,
$\alpha_h>0$, and an unpenalized intercept, suppose that the feature
map satisfies $\|\psi(x)\|_2\leq R$ on its input domain.
For any scalar $\gamma_h$, the fitted coefficients and correction satisfy
\begin{equation}
 \|w_h\|_2\le\frac{\sigma_{e,h}}{2\sqrt{\alpha_h}},\qquad
 |\gamma_h q_h(x)|\le |\gamma_h|\left[
 |\bar e_h|+(R+\|\bar\psi\|_2)
 \frac{\sigma_{e,h}}{2\sqrt{\alpha_h}}\right].
 \label{eq:mechanism-main-correction-bound}
\end{equation}
\end{proposition}
For a concatenated bounded common and departure vector,
$R=\sqrt{R_0^2+R_d^2}$; for the common-only vector, $R=R_0$.
For this decoder, the free intercept contributes the residual-mean term
and cannot be omitted. This bound controls the correction, not the baseline, total
forecast error, or sensitivity of the preceding reconstruction to changes
in the history.


\paragraph{Proof idea.}
Eliminating the intercept gives
$b_h=\bar e_h-\bar\psi^\top w_h$. If $\Psi_c$ is the matrix whose
rows are $(\psi(x_i)-\bar\psi)^\top$ and $e_c$ is the centered
residual vector, then
$w_h=(\Psi_c^\top\Psi_c+n_A\alpha_h I)^{-1}\Psi_c^\top e_c$.
The coefficient operator has singular values
$s/(s^2+n_A\alpha_h)\leq1/(2\sqrt{n_A\alpha_h})$.
This yields the coefficient bound. Applying Cauchy--Schwarz to
$q_h(x)=\bar e_h+(\psi(x)-\bar\psi)^\top w_h$ yields the
correction bound. Appendix~\ref{app:reconstruction-mechanism} gives
the complete argument, including the nonexpansiveness of the radial map.

The proposition applies separately to each selected horizon and is a property
of bounded-feature ridge regression. A summed-squared-error objective
would need penalty $n_A\alpha_h$ to define the same estimator. The
radial transformation is nonexpansive in its own input, but this does
not give a Lipschitz bound for the preceding signed reconstruction or
for its fallback boundary.

\paragraph{Baseline recovery and predictive gain.}
The calibration grid includes $\gamma_h=0$, which recovers the same
fitted baseline exactly. This is an available forecast rather than a
guarantee that a finite calibration set selects the best future forecast.
The coefficient and correction bounds likewise make no claim about
prediction accuracy. The experiment therefore evaluates error, numerical
coverage, and the correction inequalities as separate outcomes.

\subsection{Centered corrections to a frozen common predictor}
\label{sec:fixed-common-decoder}

The residual-target study fixes a separately fitted common forecast
$f_{\mathrm{com},h}$ and a shared bias before comparing departures:
\begin{equation}
 f_{\mathrm{base},h}(x)=f_{\mathrm{com},h}(x)+\beta_h,\qquad
 \beta_h=\frac1{n_B}\sum_{i\in B}
       \{y_{i,h}-f_{\mathrm{com},h}(x_i)\}.
 \label{eq:mechanism-fixed-common-baseline}
\end{equation}
The common forecast is fitted on all of A and selected on B. Its
coefficients and the B-fitted bias are then shared across the departure
comparisons and frozen before C. Thus the common component cannot
compensate for a change of departure coordinates during residual fitting.

Let $\varphi_m(x)$ be the separately standardized and radially bounded
departure block for method $m$, with $\|\varphi_m(x)\|_2\leq R_{d,m}$,
and let $\bar\varphi_m$ be its A mean. For either in-sample or
out-of-fold residual targets $r_{ih}$, write
$\bar r_h=n_A^{-1}\sum_i r_{ih}$ and fit
\begin{equation}
 v_{m,h}=\mathop{\arg\min}_{v}\left[
 \frac1{n_A}\sum_i
 \{r_{ih}-\bar r_h-(\varphi_m(x_i)-\bar\varphi_m)^\top v\}^2
 +\alpha_{m,h}\|v\|_2^2\right],\qquad \alpha_{m,h}>0.
 \label{eq:mechanism-centered-slope-fit}
\end{equation}
In-sample targets use $r_{ih}=y_{i,h}-f_{\mathrm{com},h}(x_i)$;
out-of-fold targets use the common forecast fitted without the fold
containing the episode of $x_i$. The deployed forecast is
\begin{equation}
 \widehat f_{m,h}(x)=f_{\mathrm{base},h}(x)+\eta_{m,h}\delta_{m,h}(x),
 \qquad
 \delta_{m,h}(x)=(\varphi_m(x)-\bar\varphi_m)^\top v_{m,h}.
 \label{eq:mechanism-centered-slope-forecast}
\end{equation}
It contains no additional residual-mean intercept. Subtracting the
constant shared bias before centering the fitting targets would leave
the slope unchanged. The correction has zero empirical mean on A;
its mean on another distribution need not vanish. Setting $\eta_{m,h}=0$
recovers the same shared baseline for every departure method.

\paragraph{Bound for the deployed slope.}
Define $\sigma_{r,h}^2=n_A^{-1}\sum_i(r_{ih}-\bar r_h)^2$.
The spectral argument of
Proposition~\ref{prop:mechanism-correction-bound} immediately gives
\begin{equation}
 \|v_{m,h}\|_2\leq\frac{\sigma_{r,h}}{2\sqrt{\alpha_{m,h}}},\qquad
 |\eta_{m,h}\delta_{m,h}(x)|\leq
 |\eta_{m,h}|(R_{d,m}+\|\bar\varphi_m\|_2)
 \frac{\sigma_{r,h}}{2\sqrt{\alpha_{m,h}}}.
 \label{eq:mechanism-centered-slope-bound}
\end{equation}
The residual-mean term is absent because it is absent from the deployed
correction. The fixed common forecast and shared bias remain outside
this bound. This specialization is generic bounded-feature ridge
algebra and applies equally to the matched RBF departure.

\paragraph{Residual alignment and predictive gain.}
Condition on all fitted and selected quantities, and take expectations
under a specified evaluation distribution for $(X,Y)$. Suppress the
horizon and method indices, and suppose $Y$, $f_{\mathrm{base}}(X)$,
and $\delta(X)$ have finite second moments. For fixed $\eta$, define
$\mathcal R(f)=\mathbb E[(Y-f(X))^2]$. Expansion of the square gives
the elementary identity
\begin{equation}
 \mathcal R(f_{\mathrm{base}}+\eta\delta)
 -\mathcal R(f_{\mathrm{base}})
 =\eta^2\mathbb E[\delta(X)^2]
 -2\eta\mathbb E[(Y-f_{\mathrm{base}}(X))\delta(X)].
 \label{eq:mechanism-residual-risk}
\end{equation}
For $\eta>0$, improvement holds exactly when
$\mathbb E[(Y-f_{\mathrm{base}}(X))\delta(X)]>
\eta\mathbb E[\delta(X)^2]/2$. A bound on correction
magnitude does not establish that alignment. The identity concerns
errors of the deployed full-A baseline even when fitting uses fold-model
residuals; it does not equate those two residual distributions or
identify the cause of an observed performance difference.
